\documentclass[11pt]{article}
\usepackage[a4paper,margin=1in]{geometry}
\usepackage{amsmath,amssymb,amsthm}
\usepackage[T1]{fontenc}
\usepackage{lmodern}
\usepackage[hidelinks]{hyperref}
\setlength{\emergencystretch}{3em}

\theoremstyle{plain}
\newtheorem{theorem}{Theorem}
\newtheorem{proposition}[theorem]{Proposition}
\newtheorem{lemma}[theorem]{Lemma}
\theoremstyle{definition}
\newtheorem{definition}[theorem]{Definition}
\theoremstyle{remark}
\newtheorem{remark}[theorem]{Remark}

\newcommand{\Z}{\mathbb Z}

\title{A Counterexample to Conjecture 00000008176:\\
the Artin Density of a Perfect Square}
\author{%
  Feiyu\thanks{Independent researcher.}
  \and
  Claude (Opus 5.5)\thanks{Anthropic.}%
}
\date{2026-10-03}

\begin{document}
\maketitle

\begin{abstract}
Let $\delta_A(q)$ be the density of the primes $p$ modulo which $q$ is a primitive
root. The second claim of Conjecture 00000008176 is that $\delta_A(q)$ is zero or
nonzero according to whether $q\not\equiv1\pmod{p^2}$ holds for every prime $p$. We
show unconditionally that $\delta_A(s^2)=0$ for every integer $s$. Since $q=4$
satisfies the criterion and $q=9$ violates it, the criterion decides nothing, in
either direction. The refutation is checked in Lean~4 against Mathlib.
\end{abstract}

\section{The conjecture as stated}

The original text of conjecture 00000008176 reads:

\begin{quote}
Definition: The density table for computational verification of Artin's primitive
root conjecture: the explicit values of the density delta\_A(q) of primes q being a
primitive root (Hooley explicit). Conjecture: delta\_A(2) = 0.3739558136\dots{} (the
explicit 2-value of the Artin constant, the product of $(1 - 1/(p(p-1)))$ with the
Legendre correction); the complete classification of the values of delta\_A(q) (zero
versus nonzero) is decided by q not congruent to 1 mod $p^2$ for all p (a
non-Wieferich criterion); the minimal positive value of delta is attained at q = 2 (a
2-minimality law); and the arithmeticity of the density (rationality of delta):
delta\_A(q) is rational if and only if q is a perfect square or $-4$ times a square (a
rationality criterion), the list of rational values being finite (a finite rational
list).
\end{quote}

The Chinese version in \texttt{conjectures/00000008176.md} states the same four
claims.

\section{Reading of the statement}

The conjecture is a conjunction of four claims. We refute the second one, so the
conjecture is false whatever the other three claims mean.

\paragraph{The density.} As in Artin's primitive root conjecture, $\delta_A(q)$ is
the density of the set
\[
\mathcal A(q)=\{p \text{ prime}: q \text{ is a primitive root modulo } p\}
\]
relative to the primes. Here $q$ is a primitive root modulo $p$ if its residue
generates the cyclic group $(\Z/p\Z)^\times$, that is, it has multiplicative order
$p-1$; a residue $0$ is never a primitive root. The phrase ``the density of primes q
being a primitive root'' can only mean this, since $\delta_A$ is indexed by $q$. The
fourth claim speaks of $q$ being a perfect square or $-4$ times a square, so $q$
ranges over the integers, including perfect squares.

\begin{definition}
A set $S$ of primes has \emph{density $d$} if
$\#\{p\in S: p\le x\}/\pi(x)\to d$ as $x\to\infty$, where $\pi(x)$ is the number of
primes up to $x$. We write $\delta_A(q)=0$ if $\mathcal A(q)$ has density $0$.
\end{definition}

\paragraph{The criterion.} Let $W(q)$ be the statement ``$q\not\equiv1\pmod{p^2}$
for every prime $p$''. The name ``non-Wieferich criterion'' indicates that $p$ runs
over primes. The second claim says that $W$ decides whether $\delta_A(q)$ is zero or
nonzero. The natural reading is
\begin{equation}\label{eq:N}
\delta_A(q)\ne0 \iff W(q)\qquad\text{for every integer } q. \tag{N}
\end{equation}
We also refute the reversed matching
\begin{equation}\label{eq:Z}
\delta_A(q)=0 \iff W(q)\qquad\text{for every integer } q. \tag{Z}
\end{equation}
For (N) it does not matter whether ``$\delta_A(q)\ne0$'' is read as ``$\mathcal
A(q)$ does not have density $0$'' or as ``$\mathcal A(q)$ has a density $d\ne0$''.

\section{Results}

\begin{proposition}\label{prop:square}
For every integer $s$, $\mathcal A(s^2)\subseteq\{2\}$. Consequently
$\delta_A(s^2)=0$.
\end{proposition}

\begin{theorem}\label{thm:main}
$W(4)$ holds and $\delta_A(4)=0$, so (N) fails. $W(9)$ fails and $\delta_A(9)=0$, so
(Z) fails. In particular Conjecture 00000008176 is false.
\end{theorem}

\section{Proofs}

\begin{proof}[Proof of Proposition~\ref{prop:square}]
Let $p$ be an odd prime and $a$ the residue of $s$ modulo $p$. If $a=0$, the residue
of $s^2$ is $0$, which is not a primitive root. Otherwise Fermat's little theorem
gives
\[
(a^2)^{(p-1)/2}=a^{p-1}=1,
\]
so the order of $a^2$ divides $(p-1)/2$. Since $0<(p-1)/2<p-1$, this order is not
$p-1$, and $s^2$ is not a primitive root modulo $p$. Hence $\mathcal A(s^2)
\subseteq\{2\}$.

It follows that $\#\{p\in\mathcal A(s^2):p\le x\}\le1$, so
\[
0\le\frac{\#\{p\in\mathcal A(s^2):p\le x\}}{\pi(x)}\le\frac1{\pi(x)}\to0,
\]
because $\pi(x)\to\infty$. So $\mathcal A(s^2)$ has density $0$.
\end{proof}

\begin{proof}[Proof of Theorem~\ref{thm:main}]
If $4\equiv1\pmod{p^2}$ for a prime $p$, then $p^2\mid3$, which is impossible as
$p^2\ge4$. So $W(4)$ holds, and $\delta_A(4)=0$ by
Proposition~\ref{prop:square} with $s=2$. This contradicts (N) at $q=4$ in either
reading of ``$\delta_A(4)\ne0$'': the density is $0$, and a density is unique since
it is a limit.

On the other hand $9\equiv1\pmod{2^2}$, so $W(9)$ fails, while $\delta_A(9)=0$ by
Proposition~\ref{prop:square} with $s=3$. This contradicts (Z) at $q=9$.
\end{proof}

\section{Formalization}
\sloppy

The Lean~4 project in \texttt{lean/} (file \texttt{lean/Submission/Basic.lean})
formalizes the definitions above and proves both results against Mathlib.
\begin{itemize}
  \item \texttt{artinSet q} is $\mathcal A(q)$: the primes $p$ with
    \texttt{IsPrimitiveRoot ((q : ZMod p)) (p - 1)}, using Mathlib's notion of a
    primitive root of unity, i.e.\ the residue has order exactly $p-1$.
  \item \texttt{HasPrimeDensity S d} states that
    \texttt{(S $\cap$ [0,x]).ncard / Nat.primeCounting x} tends to $d$.
  \item \texttt{NonWieferich q} is $W(q)$, with $q\equiv1\pmod{p^2}$ expressed by
    Mathlib's \texttt{Int.ModEq}.
  \item \texttt{ClassificationNonzero} is (N) with ``$\delta_A(q)\ne0$'' read as
    ``not density $0$''; \texttt{ClassificationNonzero'} is (N) with ``has a
    density $d\ne0$''; \texttt{ClassificationZero} is (Z).
  \item \texttt{artinSet\_sq\_subset} and \texttt{hasPrimeDensity\_sq} are
    Proposition~\ref{prop:square}. They use Mathlib's Fermat little theorem
    \texttt{ZMod.pow\_card\_sub\_one\_eq\_one} and
    \texttt{Nat.tendsto\_primeCounting}.
  \item \texttt{nonWieferich\_four}, \texttt{not\_nonWieferich\_nine}, and the
    refutations \texttt{not\_classificationNonzero},
    \texttt{not\_classificationNonzero'}, \texttt{not\_classificationZero}.
  \item \texttt{conjecture\_00000008176\_false} states
    $\lnot(P\wedge\texttt{ClassificationNonzero}\wedge Q\wedge R)$ for arbitrary
    propositions $P,Q,R$ standing for the other three claims.
\end{itemize}
The toolchain is \texttt{leanprover/lean4:v4.33.1}, Mathlib is pinned to
\texttt{v4.33.1}, and \texttt{lake build} succeeds. The project contains no
\texttt{sorry}, no \texttt{admit} and no \texttt{native\_decide}, and introduces no
axioms: \texttt{\#print axioms} reports exactly \texttt{propext},
\texttt{Classical.choice} and \texttt{Quot.sound} for the three refutations.

\fussy
\section{Remarks}

\begin{remark}
The refutation is unconditional: it only uses that squares have density $0$. The
expected classification is different from the stated one. If $q=-1$ or $q$ is a
perfect square, then $\delta_A(q)=0$ (the integer $-1$ is a primitive root only
modulo $2$ and $3$). Otherwise, under the generalized Riemann hypothesis, Hooley
\cite{hooley} proved that $\delta_A(q)$ exists and is positive. The criterion $W$ is
unrelated to this dichotomy in both directions: $4$ satisfies $W$ with
$\delta_A(4)=0$, and $10\equiv1\pmod{3^2}$ violates $W$ although $\delta_A(10)>0$
under the generalized Riemann hypothesis.
\end{remark}

\begin{remark}
We do not address the other three claims. Under the generalized Riemann hypothesis,
Hooley's theorem gives $\delta_A(2)$ as Artin's constant $0.3739558136\ldots$.
Unconditionally, it is not known for any single integer $q$ that $q$ is a primitive
root modulo infinitely many primes.
\end{remark}

\begin{thebibliography}{9}
\bibitem{hooley} C.~Hooley, \emph{On Artin's conjecture}, J. Reine Angew. Math. 225
  (1967), 209--220.
\end{thebibliography}

\end{document}
